\begin{abstract}
Humans learn an unfamiliar physical world by acting in it, and then see new scenes in the light of what they learned. We ask whether multimodal agents can learn the physical world through interaction in the same way. We introduce a benchmark built around a single unit, the visual experiment: a rendered scene, an intervention, and a time interval, followed by the rendered outcome. During interaction, the agent chooses experiments in a world whose laws and constants may depart from canonical physics and observes their outcomes; at fixed checkpoints, it must first predict outcomes of host-posed experiments, which traces an online learning curve. After interaction is frozen, the roles reverse: the world poses new experiments on unseen scenes, and the agent must choose the correct future frame among four rendered candidates. Worlds are organized in a five-level hierarchy that separates phenomenon families, shared constants, structural law changes, hidden coefficient fields, and individual scenes. Learning is measured as what interaction changes: the accuracy gain over no interaction, and the gap between evidence from the agent's own world and from a sibling world with a different hidden law. A measurement probe and law-informed ceilings locate where learning breaks. We evaluate frontier multimodal agents as complete model--harness systems. \TODO{Add the one-sentence quantitative headline once baseline runs finish.}
\end{abstract}
